% Detailed proofs retained separately from the submission manuscript.
% Deliberately not input by main.tex; the main text contains proof sketches.
% Keep this file for the full technical version only.
\appendices
\section{Proofs of Formal Results}
\label{app:proofs}
We use $\Sig\models\gamma$ in its standard rule-consequence sense:
every model of $\Sig$ satisfies every ground instance of $\gamma$.

\begin{proof}[Proof of Lemma~\ref{lem:closurepreservation}]
Consider one Phase~1 deletion step. Let $G'\subseteq G$ be the current
graph and $\gamma$ the non-recursive rule used in this step.
The subgraph deleted in this step is $\Gdel=G'\cap\tau_{\{\gamma\},G'}$. Every fact in $\Gdel$
has the head predicate of $\gamma$, which occurs in no body atom of
$\gamma$. Consequently, for each $d\in \Gdel$, any body match in $G'$
witnessing $d$ remains a body match in $G'\setminus \Gdel$. Thus
$\Gdel\subseteq\Delta_{G'\setminus \Gdel,\Sig}$ and
$G'\subseteq\Delta_{G'\setminus \Gdel,\Sig}$. Since $G'\setminus \Gdel\subseteq G'$,
monotonicity and idempotence of least-fixpoint closure imply
\[
\Delta_{G'\setminus \Gdel,\Sig}=\Delta_{G',\Sig}.
\]
Applying this argument to every Phase~1 deletion proves that its
output $G'$ is a subset of $G$ and has closure $\Delta_{G,\Sig}$.
The argument does not depend on the particular SCC or rule order.

For a Phase~2 round, let $G'$ be its input graph and let
$G'_i=\bigcup_{j=1}^{i}\Gblock^{(j)}$ be the subgraph formed by its first $i$ blocks,
with $G'_0=\emptyset$.
Denote the retained set and its closure after block $i$ by $\Gkeep^{(i)}$
and $\Gcl^{(i)}=\Delta_{\Gkeep^{(i)},\Sig}$, with $\Gkeep^{(0)}=\Gcl^{(0)}=\emptyset$.
The initial closure is empty because PD rules have nonempty bodies.
We claim, by induction on $i$, that
\[
\Gkeep^{(i)}\subseteq G'_i\subseteq \Gcl^{(i)}.
\]
It holds for $i=0$. At block $i$, the update
$\Gkeep^{(i)}=\Gkeep^{(i-1)}\cup(\Gblock^{(i)}\setminus \Gcl^{(i-1)})$ retains only input facts,
so $\Gkeep^{(i)}\subseteq G'_i$. Every fact of $G'_{i-1}$ lies in
$\Gcl^{(i-1)}\subseteq \Gcl^{(i)}$. A fact of $\Gblock^{(i)}$ either already lies in
$\Gcl^{(i-1)}$ or is added to $\Gkeep^{(i)}\subseteq \Gcl^{(i)}$. Hence
$G'_i\subseteq \Gcl^{(i)}$, proving the induction step.
At the last block, the partition gives $G'_{N_b}=G'$, so
$\Gkeep^{(N_b)}\subseteq G'\subseteq\Delta_{\Gkeep^{(N_b)},\Sig}$.
Monotonicity and idempotence yield
$\Delta_{G',\Sig}=\Delta_{\Gkeep^{(N_b)},\Sig}$. Each round therefore preserves its
input closure while retaining a subset of its input. Induction over
the $T$ rounds, following Phase~1, proves
$\Gc\subseteq G$ and $\Delta_{\Gc,\Sig}=\Delta_{G,\Sig}$.
All loops terminate because the graphs and rule set are finite and
the function-free closure reaches a finite fixpoint.
\end{proof}

\begin{proof}[Proof of Theorem~\ref{thm:cot}]
We first justify the derivation binding returned by
\algref{alg:derive}. Define $G^{(0)}=\Gc$ and
$G^{(i)}=G^{(i-1)}\cup\tau_{\Sig,G^{(i-1)}}$ for $i\geq1$.
At the start of iteration $i$, the algorithm has $\Gknown=G^{(i-1)}$
and, for $i>1$, $\Gfront=G^{(i-1)}\setminus G^{(i-2)}$;
for $i=1$, $\Gfront=G^{(0)}$. Every head admitted by the algorithm
belongs to $G^{(i)}\setminus G^{(i-1)}$, because all of its body facts
belong to $G^{(i-1)}$ and previously known heads are excluded.
Conversely, let $d\in G^{(i)}\setminus G^{(i-1)}$ and choose a rule
instance deriving it from $G^{(i-1)}$. At least one body occurrence
belongs to $\Gfront$. Otherwise all its body facts would
already belong to $G^{(i-2)}$, and $d$ would be in $G^{(i-1)}$.
For $i=1$, every body occurrence belongs to the initial $\Gfront$.
The algorithm examines this position and therefore admits $d$.
It records one witness for each newly admitted head, even if several
instances derive the same head. Induction proves that the returned
$G^{*}$ is $\Delta_{\Gc,\Sig}$ and that every
$d\in G^{*}\setminus\Gc$ has a recorded witness whose body facts
have strictly smaller derivation levels. Finiteness of the active
domain makes this process terminate.

Now consider $\FunctionName{UnfoldRule}$ for an omitted fact $d$.
Maintain the following invariant for its current rule $\gamma^{(j)}$ and
ground witness $\sigma_{\mathrm w}$:
\[
\Sig\models\gamma^{(j)},\qquad
\sigma_{\mathrm w}(H_{\gamma^{(j)}})=d,\qquad
\sigma_{\mathrm w}(B_{\gamma^{(j)}})\subseteq G^{*}.
\]
It holds initially for the freshly renamed input rule recorded for
$d$. If a body occurrence $\beta$ has witness
$e=\sigma_{\mathrm w}(\beta)\notin\Gc$, the binding $\Der[e]$
provides a freshly renamed rule $\gamma'$ and ground witness
$\sigma'$ with $\sigma'(H_{\gamma'})=e$.
The two clauses have disjoint variables. Their ground witnesses
agree on the resolved atom, so their union unifies
$\beta$ and $H_{\gamma'}$. The relevant MGU $\mu$ therefore
factors this ground unifier on the resolved atoms. Extending the
factorization by the original ground assignments on other clause
variables gives the ground substitution $\lambda$ used in the
algorithm.
Applying $\mu$ to the whole rule and replacing just the selected
body occurrence by $\mu(B_{\gamma'})$ preserves the witnessed
head and assigns the child's recorded body facts to the inserted
occurrences. Definite-clause resolution also preserves
$\Sig\models\gamma^{(j+1)}$. Range restriction is preserved: a head variable
whose only body occurrence was replaced remains represented in the
range-restricted child's body. Thus the invariant continues to hold.

Each expanded occurrence has a derivation level greater than the
levels of its inserted supports. Expansion therefore forms a
finitely branching tree of depth at most $\ell$, the number of
productive derivation levels. It terminates with every ground body
fact in $\Gc$. Keeping occurrences distinct until the end ensures
that only the chosen occurrence is expanded. Removing duplicate
body atoms on return changes neither the clause's meaning nor its
ground witness. Hence the returned path rule $\gamma_d$ is
entailed by $\Sig$ and has a ground instance with head $d$ and
body contained in $\Gc$.

Every initialized identity rule is a tautology and covers its
retained facts in one step; pruning retains these identity rules.
During the loop in
\algref{alg:gencot}, each omitted fact is either already
covered by a current rule or receives such a path rule. Therefore,
before pruning, every fact in $G^{*}$ is covered by an instance
over $\Gc$. Deduplication preserves this property.
Suppose a remaining rule $\gamma'$ subsumes a rule $\gamma$
being removed, witnessed by a substitution $\theta$ with
$\theta(H_{\gamma'})=H_\gamma$ and
$\theta(B_{\gamma'})\subseteq B_\gamma$.
For any ground witness $\sigma$ of $\gamma$ over $\Gc$,
$\sigma\circ\theta$ witnesses the same head through $\gamma'$,
since
\[
(\sigma\circ\theta)(B_{\gamma'})
 \subseteq\sigma(B_\gamma)\subseteq\Gc.
\]
Each pruning step therefore preserves coverage. All surviving
rules are entailed by $\Sig$, giving
$\Gc\cup\tau_{\Sigcot,\Gc}\subseteq\Delta_{\Gc,\Sig}$;
coverage gives the reverse inclusion. Thus $\Sigcot$ satisfies all
conditions of Definition~\ref{def:cot}. At most one finite path
rule is generated per omitted closure fact, in addition to finitely
many identities, so $\Sigcot$ is finite. None of these operations
uses the query.
\end{proof}

\begin{proof}[Proof of Theorem~\ref{thm:queryeq}]
Lemma~\ref{lem:closurepreservation} gives
$G^{*}=\Delta_{\Gc,\Sig}$, and Theorem~\ref{thm:cot} gives
$G^{*}=\Gc\cup\tau_{\Sigcot,\Gc}$.
The initialized identities cover every fact of $\Gc$, and the
pruning argument in the preceding proof preserves that coverage.
Consequently, the stronger atomic equality
\[
\tau_{\Sigcot,\Gc}=G^{*}
\]
holds. This equality is useful because atom-wise rewriting uses
rule candidates, rather than a separate stored-fact branch.

Fix a query atom $\alpha_k$ with distinct variables $\vect{Z}_k$.
Let $R_k^{*}$ be its match relation over $G^{*}$, with columns
$\vect{Z}_k$. For any candidate rule, a body match over $\Gc$
produces an instance of its head in $G^{*}$ by soundness. Unifying
that head with $\alpha_k$ therefore produces a tuple in $R_k^{*}$;
the output mapping preserves constants and repeated variables.
Conversely, take any tuple in $R_k^{*}$ and its ground match of
$\alpha_k$. By the atomic equality, a ground instance of some
$\gamma\in\Sigcot$ has the same head fact and a body match in $\Gc$.
Rename $\gamma$ apart from the query. The two head matches give
a ground unifier of $H_\gamma$ and $\alpha_k$.
Relevant-MGU factorization yields a match of the corresponding
candidate body, extending the ground substitution to body-only
rule variables using the original rule-instance witness. Its
output tuple is precisely the selected tuple
of $R_k^{*}$. Hence the union of candidate branches equals $R_k^{*}$.

The query frontier $\xi_k$ contains every variable of $\alpha_k$
that either occurs in another query atom or is an answer variable.
All other variables of $\alpha_k$ can be existentially quantified
within this group. Therefore its compiled relation is
$R_k=\pi_{\xi_k}(R_k^{*})$. By ordinary CQ semantics, the answer relation
for $Q$ over $G^{*}$ is
\[
\pi_{\vect{X}}\bigl(R_1\bowtie\cdots\bowtie R_m\bigr).
\]
This also holds for disconnected atoms, constants, and identified
variables. With an empty frontier, $R_k$ is either the unit
zero-column tuple or the empty relation, so it contributes exactly
the required existence condition. The same convention handles
Boolean queries. This proves P00.

For \QFSE{}, fix any ordering $\beta_1,\ldots,\beta_s$ of a
candidate body's atom occurrences. The retained variable set
$\xi_{k,i}^{(j)}$ in~\S\ref{sec:sqlcomposition} contains exactly the variables
among the first $j$ atoms needed by the output or a later atom.
The stage invariant is
\[
R_{k,i}^{(j)}
=\pi_{\xi_{k,i}^{(j)}}
 (R_{\beta_1}\bowtie\cdots\bowtie R_{\beta_j}),
\]
with
$R_{k,i}^{(0)}=\{()\}$. A variable projected out at stage $j-1$
occurs neither in the branch's required output nor in any
remaining atom; hence it cannot influence a later join.
Joining $R_{k,i}^{(j-1)}$ with $R_{\beta_j}$ and projecting onto
$\xi_{k,i}^{(j)}$
establishes the invariant by induction. If every variable in
$\xi_{k,i}^{(j)}$ is already a column of $R_{k,i}^{(j-1)}$, the current atom
introduces no required output variable. Under set semantics,
\[
\pi_{\xi_{k,i}^{(j)}}(R_{k,i}^{(j-1)}\bowtie R_{\beta_j})
=\pi_{\xi_{k,i}^{(j)}}(R_{k,i}^{(j-1)}\ltimes R_{\beta_j}),
\]
so the permitted semi-join implements the same stage. Constant
and repeated-variable constraints belong to each atom's match
relation. At the final stage, the retained variables and any
constant output terms restore exactly the original branch's
$\xi_k$ columns. Thus \QFSE{} preserves every group, regardless of
the chosen body-atom order.

For \UFA{}, each source group with $\xi_k=\{x\}$ is a unary relation
on $x$. Their natural join is exactly
$R_x^{\mathrm u}=\bigcap_{k\in \mathcal I_x}\pi_x(R_k)$.
For a multi-frontier group $R_j$ containing $x$, replacing it by
$R_j\ltimes_x R_x^{\mathrm u}$ pushes the same membership constraint into
that group. If at least one such group remains, its $x$ column
keeps the constraint in the top-level join, so the original
singleton-frontier sources can be removed. Applying the filter
to every eligible group is idempotent under set semantics.
If no multi-frontier group contains $x$, retaining $R_x^{\mathrm u}$ when
$x$ is an answer variable preserves its possible output values;
when $x$ is existential, an $\mathrm{EXISTS}(R_x^{\mathrm u})$ guard
preserves its required nonemptiness. Every zero-frontier group
likewise remains an existence guard. These rules also cover an
empty top-level join.

In P11, the $R_x^{\mathrm u}$ membership test is applied to each candidate
using its mapped output value, after that value is bound or
immediately if it is a constant. A frontier variable needed by
the test is included in the fixed required output of \QFSE{} and
cannot be projected away before or after the test. Thus the
\QFSE{} and \UFA{} equivalences compose.
\algref{alg:compile} emits the unions, joins, projections,
semi-joins, and existence guards just proved equivalent, with
distinct outputs under the paper's set semantics. Its SQL
interpretation uses a faithful null-free encoding of $\Gc$ and
equality joins. For completed executions, each of P00, P01, P10,
and P11 therefore returns exactly the answers to $Q$ over $G^{*}$.
Selecting one of these equivalent modes with the cost-sensitive
policy tree cannot change the answer set. Taking a finite union
of CQs with a common answer schema also preserves the equality
branch by branch.
\end{proof}

\subsection{Auxiliary Size and Cost Bounds}
For completeness, the size bound in~\eqref{eq:cot-size-bound}
follows from the same decreasing-level argument. Each productive
iteration of \algref{alg:derive} creates at least one of
the $M=|G^{*}\setminus\Gc|$ omitted facts, so $\ell\leq M$.
A path rule's unfolding tree has branching at most
$b=\max(\{1\}\cup\{|B_\gamma|:\gamma\in\Sig\})$
and depth at most $\ell$. It has at most $b^{\ell}$ leaf
occurrences, hence at most that many body atoms after
deduplication. There are at most $M$ generated path rules, each
with one head, and $P$ identities, each with one head and one
body atom. Thus
\[
L_{\Sigcot}\leq 2P+M(1+b^{\ell}).
\]
Deduplication and subsumption removal can only decrease this
count. The estimate establishes finiteness; it is exponential
in $\ell$ when $b>1$.

The grouped candidate representation stores one group record
per query atom, one record per candidate, and each candidate
body once. With unary/binary predicates and interned symbols,
the output mapping has constant size per candidate. For $m$
atoms, $n_k$ candidates in group $k$, and total candidate-body
length $L$, this gives $O(m+\sum_k n_k+L)$ words, as stated in
\S\ref{sec:correctness}. The bound concerns the grouped
representation, not emitted SQL or database intermediate results.
Compiling four modes, extracting their features, traversing a
height-$h$ policy tree, and executing one selected plan gives
\eqref{eq:online-cost}. Counting the retained facts, the
head/body atoms of $\Sigcot$, and the policy model gives
\eqref{eq:stored-space} under the stated bounded-arity and
interned-symbol convention.
